\hypertarget{appendix-749ba51a6c25}{}\label{referencestart:appendix-749ba51a6c25}
\section{GPU vertex specification}\label{appendix-749ba51a6c25-gpu-vertex-specification}

\textbf{Curriculum depth:} Developed · \textbf{First seen:} University courses · \textbf{Latest development:} University courses

\begin{allmightroledefinition}

\subsection{Definition and mechanism}\label{appendix-749ba51a6c25-definition-and-mechanism}

A GPU vertex-input specification explains how bytes in one or more buffers become per-vertex shader inputs. A buffer object owns bytes; a vertex array object records their interpretation: enabled attributes, component count and type, normalization, stride, byte offset, and the buffer binding used for fetch. Stride is the number of bytes between successive vertices; offset is the first component's byte position. For vertex i, the address is base + offset + i × stride. For interleaved records containing position (3 floats), normal (3), and UV (2), the stride is 8 × 4 = 32 bytes and the UV offset is 24 bytes. Modern binding APIs separate attribute format from the buffer binding; older pointer-style calls can capture the buffer bound at setup time. A wrong stride may still address valid memory, so the symptom is often plausible but scrambled geometry. The contract must agree with CPU packing, shader input types, and the API's state-capture rules.

\end{allmightroledefinition}

\subsection{Prerequisites}\label{appendix-749ba51a6c25-prerequisites}

\begin{itemize}
\tightlist
\item
  \href{../07-Vision-XR-and-Playable-Worlds/main.pdf\#nameddest=appendix-26bf6bf2fb8b}{Vector storage SIMD and boundary layouts (Vision, XR, and Playable Worlds, Reference notes: Vector storage SIMD and boundary layouts)}
\item
  \href{../02-Mathematics-for-Visual-Computing/main.pdf\#nameddest=appendix-ffdc1d4e5350}{Matrix storage indexing and strides (Mathematics for Visual Computing, Reference notes: Matrix storage indexing and strides)}
\item
  \hyperref[appendix-1c07d64f5e70]{Camera view and clip-space pipeline}
\item
  \href{../01-Bits-Programs-and-Reliable-Systems/main.pdf\#nameddest=appendix-6af5921bc853}{Pointer arithmetic and provenance (Bits, Programs, and Reliable Systems, Reference notes: Pointer arithmetic and provenance)}
\end{itemize}

\subsection{Used by}\label{appendix-749ba51a6c25-used-by}

\begin{itemize}
\tightlist
\item
  \href{../02-Mathematics-for-Visual-Computing/main.pdf\#nameddest=03-visual-computing-i-ch012-rasterization-coverage-and-depth}{Rasterization, coverage, and depth (Mathematics for Visual Computing, Ch. 16)}
\item
  \href{../03-Geometry-Cameras-and-Image-Formation/main.pdf\#nameddest=appendix-40426e4b2103}{Homogeneous clipping and perspective interpolation (Geometry, Cameras, and Image Formation, Reference notes: Homogeneous clipping and perspective interpolation)}
\end{itemize}

\subsection{Full teaching explanation}\label{appendix-749ba51a6c25-full-teaching-explanation}

\begin{itemize}
\tightlist
\item
  \href{../02-Mathematics-for-Visual-Computing/main.pdf\#nameddest=03-visual-computing-i-ch012-rasterization-coverage-and-depth}{Rasterization, coverage, and depth (Mathematics for Visual Computing, Ch. 16)}
\end{itemize}

\subsection{Recall and next depth}\label{appendix-749ba51a6c25-recall-and-next-depth}

\begin{itemize}
\tightlist
\item
  Trace the fetched address for vertex 3 when stride is 32 bytes and offset is 24 bytes.
\item
  Why can an incorrect stride produce plausible but scrambled geometry without a memory fault?
\end{itemize}

\subsection{Source and navigation}\label{appendix-749ba51a6c25-source-and-navigation}

Teaching narrative: \href{../02-Mathematics-for-Visual-Computing/main.pdf\#nameddest=03-visual-computing-i-ch012-rasterization-coverage-and-depth}{open the main explanation (Mathematics for Visual Computing, Ch. 16)}. Browse the \hyperref[appendix-a6d556a4ef50-university-graphics-vision-and-photography---independent-series]{University graphics and vision index}. Additional chapter connections: \href{../03-Geometry-Cameras-and-Image-Formation/main.pdf\#nameddest=03-visual-computing-i-ch001-a-scene-description-becomes-an-image}{A scene description becomes an image (Geometry, Cameras, and Image Formation, Ch. 8)}; \href{../02-Mathematics-for-Visual-Computing/main.pdf\#nameddest=03-visual-computing-i-ch012-rasterization-coverage-and-depth}{Rasterization, coverage, and depth (Mathematics for Visual Computing, Ch. 16)}; \href{../03-Geometry-Cameras-and-Image-Formation/main.pdf\#nameddest=02-foundations-ii-ch038-the-mathematical-table-needs-an-explicit-storage-agreement}{The mathematical table needs an explicit storage agreement (Geometry, Cameras, and Image Formation, Ch. 7)}; \href{../01-Bits-Programs-and-Reliable-Systems/main.pdf\#nameddest=01-foundations-i-ch003-a-pointer-carries-an-access-contract}{A pointer carries an access contract (Bits, Programs, and Reliable Systems, Ch. 3)}.

Caption/source records: \href{../07-Vision-XR-and-Playable-Worlds/main.pdf\#nameddest=04-visual-computing-ii-ch036-p1-01---module-1.1-course-intro-administrivia-part-1-graphics-visual-computing-2026}{P1-01 (Vision, XR, and Playable Worlds, Ch. 5)}, \href{../07-Vision-XR-and-Playable-Worlds/main.pdf\#nameddest=04-visual-computing-ii-ch036-p3-03---moving-objects-in-space}{P3-03 (Vision, XR, and Playable Worlds, Ch. 5)}, \href{../07-Vision-XR-and-Playable-Worlds/main.pdf\#nameddest=04-visual-computing-ii-ch036-p3-04---rotations-about-an-arbitrary-axis}{P3-04 (Vision, XR, and Playable Worlds, Ch. 5)}.
